\documentclass[11pt,a4paper]{article}

% =========================================================================
%  RAPPORT IOT — Réseaux sans fil pour l'IoT
%  Theme : IoT — palette teal / vert
% =========================================================================

\usepackage[utf8]{inputenc}
\usepackage[T1]{fontenc}
\usepackage[french]{babel}
\usepackage{lmodern}
\usepackage[a4paper,margin=2.2cm,top=2.6cm,bottom=2.4cm]{geometry}
\usepackage{xcolor}
\usepackage{fontawesome5}
\usepackage{tikz}
\usetikzlibrary{shapes.geometric,positioning,calc,arrows.meta,decorations.pathreplacing}
\usepackage{titlesec}
\usepackage[most]{tcolorbox}
\usepackage{listings}
\usepackage{fancyhdr}
\usepackage{booktabs}
\usepackage{tabularx}
\usepackage{array}
\usepackage{colortbl}
\usepackage{enumitem}
\usepackage{microtype}
\usepackage[bookmarks=true,colorlinks=true,linkcolor=accent,urlcolor=accent,citecolor=accent]{hyperref}

% --- Palette IoT ---------------------------------------------------------
\definecolor{primary}{HTML}{0F2A2E}
\definecolor{accent}{HTML}{10B981}
\definecolor{accentdark}{HTML}{047857}
\definecolor{light}{HTML}{ECFDF5}
\definecolor{midgray}{HTML}{6B7280}
\definecolor{codebg}{HTML}{1B2A2E}
\definecolor{codefg}{HTML}{E6F4EA}
\definecolor{codekw}{HTML}{5EE2A0}
\definecolor{codestr}{HTML}{F5D67A}
\definecolor{codecmt}{HTML}{7AA197}

\titleformat{\section}{\sffamily\Large\bfseries\color{primary}}{\thesection}{0.8em}{}[\vspace{-0.4ex}{\color{accent}\rule{\linewidth}{1.2pt}}]
\titleformat{\subsection}{\sffamily\large\bfseries\color{accent}}{\thesubsection}{0.6em}{}
\titlespacing*{\section}{0pt}{1.6em}{0.8em}

\pagestyle{fancy}
\fancyhf{}
\renewcommand{\headrulewidth}{0pt}
\fancyhead[L]{\footnotesize\sffamily\color{midgray}\faMicrochip~IoT — Réseaux sans fil}
\fancyhead[R]{\footnotesize\sffamily\color{midgray}Grondin A. \textbullet{} Honorine K. \textbullet{} Grondin B.}
\fancyfoot[C]{\footnotesize\sffamily\color{midgray}\thepage}

\lstdefinestyle{pystyle}{
  backgroundcolor=\color{codebg},
  basicstyle=\ttfamily\footnotesize\color{codefg},
  keywordstyle=\color{codekw}\bfseries,
  stringstyle=\color{codestr},
  commentstyle=\color{codecmt}\itshape,
  showstringspaces=false, breaklines=true,
  frame=none, framesep=8pt, framerule=0pt,
  xleftmargin=12pt, xrightmargin=12pt,
  columns=fullflexible, inputencoding=utf8, extendedchars=true,
  literate={é}{{\'e}}1 {è}{{\`e}}1 {ê}{{\^e}}1 {à}{{\`a}}1 {ô}{{\^o}}1 {î}{{\^\i}}1 {ç}{{\c{c}}}1 {ù}{{\`u}}1 {â}{{\^a}}1 {É}{{\'E}}1 {È}{{\`E}}1 {Ê}{{\^E}}1 {À}{{\`A}}1 {Ô}{{\^O}}1 {Ç}{{\c{C}}}1 {Ù}{{\`U}}1 {Â}{{\^A}}1
}
\lstset{style=pystyle,language=Python}

\newtcblisting{codebox}[1][]{
  listing only,
  enhanced jigsaw, breakable,
  colback=codebg, colframe=accent,
  arc=2pt, boxrule=0pt, leftrule=3pt,
  left=8pt, right=8pt, top=4pt, bottom=4pt,
  listing options={style=pystyle,language=Python}, #1
}

\newtcolorbox{infobox}[1][]{
  enhanced, colback=light, colframe=accent,
  arc=2pt, boxrule=0pt, leftrule=3pt,
  left=10pt, right=10pt, top=6pt, bottom=6pt,
  fontupper=\small, #1
}
\newtcolorbox{archbox}[1][]{
  enhanced, colback=accent!10, colframe=accent,
  arc=2pt, boxrule=0pt, leftrule=3pt,
  left=10pt, right=10pt, top=6pt, bottom=6pt,
  fontupper=\small, #1
}

\setlist[itemize]{leftmargin=*,itemsep=2pt,topsep=4pt}
\setlist[enumerate]{leftmargin=*,itemsep=2pt,topsep=4pt}

\begin{document}
\pagenumbering{gobble}

\begin{tikzpicture}[remember picture, overlay]
  \fill[primary] (current page.north west) rectangle ([yshift=-9cm]current page.north east);
  \fill[accent] ([yshift=-9cm]current page.north west) rectangle ([yshift=-9.25cm]current page.north east);
  % Logo : microchip + waves
  \node[anchor=center] at ([yshift=-4.6cm]current page.north) {%
    \begin{tikzpicture}[scale=1.0]
      % Microchip
      \fill[white,rounded corners=3pt] (-1.1,-1.1) rectangle (1.1,1.1);
      \fill[accent,rounded corners=2pt] (-0.85,-0.85) rectangle (0.85,0.85);
      % Pins
      \foreach \y in {-0.7,-0.35,0,0.35,0.7} {
        \fill[white] (-1.4,\y-0.07) rectangle (-1.1,\y+0.07);
        \fill[white] (1.1,\y-0.07) rectangle (1.4,\y+0.07);
        \fill[white] (\y-0.07,-1.4) rectangle (\y+0.07,-1.1);
        \fill[white] (\y-0.07,1.1) rectangle (\y+0.07,1.4);
      }
      % Centre : node
      \fill[white] (0,0) circle (0.32);
      \node[scale=0.9,accent] at (0,0) {\faIcon{wifi}};
      % Ondes Wi-Fi
      \foreach \r/\o in {2.0/0.6, 2.5/0.4, 3.0/0.25} {
        \draw[white,line width=1.2pt,opacity=\o] (0,0) ++(45:\r) arc (45:135:\r);
        \draw[white,line width=1.2pt,opacity=\o] (0,0) ++(225:\r) arc (225:315:\r);
      }
    \end{tikzpicture}};
  \node[anchor=north, white, font=\sffamily\Huge\bfseries] at ([yshift=-6.7cm]current page.north) {Réseaux IoT};
  \node[anchor=north, accent, font=\sffamily\Large] at ([yshift=-7.65cm]current page.north) {FiPy \textbullet{} MQTT \textbullet{} Grafana};
  \node[anchor=north, white!75!primary, font=\sffamily\small] at ([yshift=-8.4cm]current page.north) {R4.IOM.09 — Acquisition à supervision};

  \node[anchor=south west, primary, font=\sffamily\small] at ([xshift=2.2cm,yshift=3.3cm]current page.south west) {\textbf{\textcolor{accent}{Équipe}}};
  \node[anchor=south west, primary, font=\sffamily\footnotesize] at ([xshift=2.2cm,yshift=2.85cm]current page.south west) {GRONDIN Angélique};
  \node[anchor=south west, primary, font=\sffamily\footnotesize] at ([xshift=2.2cm,yshift=2.55cm]current page.south west) {HONORINE Kylian};
  \node[anchor=south west, primary, font=\sffamily\footnotesize] at ([xshift=2.2cm,yshift=2.25cm]current page.south west) {GRONDIN Benjamin};
  \node[anchor=south east, primary, font=\sffamily\small] at ([xshift=-2.2cm,yshift=3cm]current page.south east) {\textbf{\textcolor{accent}{Encadrant}}};
  \node[anchor=south east, primary, font=\sffamily\small] at ([xshift=-2.2cm,yshift=2.5cm]current page.south east) {MOURAD Nour};
  \fill[accent] ([yshift=1.3cm]current page.south west) rectangle ([yshift=1.4cm]current page.south east);
  \node[anchor=south, midgray, font=\sffamily\footnotesize] at ([yshift=0.6cm]current page.south) {IUT R\&T — BUT 3 — Université de La Réunion};
\end{tikzpicture}
\newpage

\pagenumbering{arabic}
\setcounter{page}{1}

{\sffamily
\hypersetup{linkcolor=primary}
\renewcommand{\contentsname}{\textcolor{accent}{\faList~} Sommaire}
\tableofcontents
}
\vspace{0.4cm}

\section{Objectifs}

\begin{infobox}[title={\faBullseye~Mission}]
Déployer une chaîne IoT complète \textemdash{} de la mesure physique à la visualisation \textemdash{} mobilisant un microcontrôleur Pycom FiPy, des capteurs embarqués PySense, le protocole MQTT, et un dashboard Grafana temps réel.
\end{infobox}

Quatre objectifs structurent ce TP :
\begin{enumerate}
  \item Configurer un réseau sans fil local FiPy $\rightarrow$ broker MQTT
  \item Acquérir les mesures des capteurs et les sérialiser en JSON
  \item Afficher les données en temps réel sur Grafana
  \item Intégrer un retour visuel via la LED RGB embarquée
\end{enumerate}

\section{Architecture du système}

\begin{archbox}[title={\faProjectDiagram~Chaîne de traitement}]
\textbf{Capteurs PySense} $\rightarrow$ \textbf{FiPy (ESP32)} $\rightarrow$ \textbf{Wi-Fi (smartphone)} $\rightarrow$ \textbf{Mosquitto (PC)} $\rightarrow$ \textbf{Grafana}
\end{archbox}

\renewcommand{\arraystretch}{1.3}
\begin{tabularx}{\linewidth}{@{}>{\bfseries}p{2.8cm} X@{}}
\toprule
\rowcolor{light} Composant & Rôle \\ \midrule
PySense       & Collecte des grandeurs physiques via I²C \\
FiPy (ESP32)  & Microcontrôleur \textemdash{} logique embarquée, conversion, publication MQTT \\
PC            & Environnement de développement + hébergement de Mosquitto \\
Smartphone    & Passerelle Wi-Fi pour contourner les restrictions réseau pédagogique \\
Mosquitto     & Broker MQTT : centralise les messages et les redirige par topics \\
Grafana       & Visualisation dynamique des séries temporelles \\
\bottomrule
\end{tabularx}

\medskip
Cette architecture en couches respecte le principe IoT fondamental : \textbf{collecter, transmettre, traiter, visualiser, agir}.

\section{Protocoles utilisés}

\subsection{USB — alimentation et débogage}
La liaison USB entre FiPy et PC remplit deux fonctions : alimentation 5 V et communication série pour les logs MicroPython via l'extension PyMakr de VS Code.

\subsection{Wi-Fi — transport TCP/IP}
Le FiPy est configuré en mode \textbf{station (STA)}. Il rejoint le point d'accès du smartphone, obtient une IP par DHCP, puis ouvre une connexion TCP vers le port \textbf{1883} (MQTT en clair) du broker. Un mécanisme de reconnexion automatique gère les pertes momentanées.

\subsection{MQTT — publication/abonnement}
MQTT est le standard de fait pour l'IoT : faible empreinte, architecture pub/sub, résilience aux pertes de connexion. Le FiPy publie sur des topics dédiés :

\begin{codebox}
/iot/groupe2gh/temperature
/iot/groupe2gh/humidite
/iot/groupe2gh/luminosite
/iot/groupe2gh/altitude
\end{codebox}

Le broker Mosquitto centralise les messages puis les redirige vers les clients abonnés (ici Grafana).

\section{Capteurs et actionneur}

\subsection{Capteurs PySense}
Tous interfacés via le bus I²C (SDA/SCL), 2 fils suffisent.

\begin{tabularx}{\linewidth}{@{}>{\bfseries}l l c c c@{}}
\toprule
\rowcolor{light} Capteur & Mesure & Type & Précision & Fréq. \\ \midrule
BME280       & T° / humidité / pression  & I²C    & $\pm 1$°C, $\pm 3$\%HR & 0,5 Hz \\
LTR329ALS01  & Luminosité                & I²C    & $\pm 5$\%               & 0,5 Hz \\
(calculée)   & Altitude (issue pression) & calc.  & $\pm 1$\,m              & 0,5 Hz \\
\bottomrule
\end{tabularx}

\subsection{Actionneur — LED RGB}
La LED intégrée du FiPy joue le rôle d'indicateur d'état :
\begin{itemize}
  \item \textcolor{blue!70!black}{\faCircle} \textbf{Bleu} \textemdash{} température inférieure à 25\,°C
  \item \textcolor{violet}{\faCircle} \textbf{Violet} \textemdash{} température supérieure à 25\,°C
  \item \textcolor{red}{\faCircle} \textbf{Rouge} \textemdash{} perte de connexion Wi-Fi
\end{itemize}

\section{Format JSON et flux de publication}

Chaque acquisition est sérialisée en JSON avant publication :

\begin{codebox}[listing options={style=pystyle,language=Java}]
{
  "temperature": 30.03,
  "humidite":    64.0,
  "luminosite":  882,
  "altitude":    10.2
}
\end{codebox}

Cinq étapes structurent la boucle principale en MicroPython :

\begin{enumerate}
  \item \textbf{Initialisation} \textemdash{} modules \texttt{network}, \texttt{pysense}, \texttt{SI7006A20}, \texttt{LTR329ALS01}, \texttt{pycom}, \texttt{umqtt}
  \item \textbf{Connexion Wi-Fi} \textemdash{} mode STA, attente DHCP
  \item \textbf{Acquisition} \textemdash{} lecture cyclique des capteurs (toutes les 2\,s)
  \item \textbf{Publication MQTT} \textemdash{} envoi sur les topics dédiés, gestion des reconnexions
  \item \textbf{Indication LED} \textemdash{} mise à jour de la couleur selon les seuils
\end{enumerate}

\section{Stockage MySQL et Grafana}

\subsection{Base de données \texttt{iot.mesures}}

Une table relationnelle persiste les acquisitions au-delà du flux MQTT :

\begin{tabularx}{\linewidth}{@{}>{\ttfamily\bfseries}l X@{}}
\toprule
\rowcolor{light} \normalfont\textbf{Colonne} & Description \\ \midrule
ts\_epoch       & Timestamp UNIX de la mesure \\
t\_c            & Température en degrés Celsius \\
h\_pct          & Humidité en pourcentage \\
lux\_ch0/ch1    & Luminosité sur deux canaux \\
alt\_m          & Altitude calculée en mètres \\
src\_device     & Identifiant émetteur (ex. \texttt{fipy1}) \\
ts\_ingest      & Horodatage lisible de l'insertion \\
\bottomrule
\end{tabularx}

\subsection{Requête type Grafana}

\begin{codebox}[listing options={style=pystyle,language=SQL}]
SELECT
  FROM_UNIXTIME(ts_epoch) AS time,
  t_c                     AS value,
  src_device              AS metric
FROM mesures
ORDER BY ts_epoch;
\end{codebox}

La fonction \texttt{FROM\_UNIXTIME} convertit le timestamp en format temporel utilisable par Grafana, qui interprète alors les données comme une série temporelle.

\subsection{Dashboards}

Quatre panels \texttt{Gauge} avec seuils colorés affichent la dernière valeur reçue :
\begin{itemize}
  \item \textcolor{green!60!black}{\textbf{Vert}} \textemdash{} valeur dans la plage normale
  \item \textcolor{orange}{\textbf{Orange}} \textemdash{} valeur de vigilance
  \item \textcolor{red}{\textbf{Rouge}} \textemdash{} valeur critique
\end{itemize}

Rafraîchissement à 5\,s : bon compromis entre réactivité et stabilité du rendu.

\section{Performances et énergie}

\renewcommand{\arraystretch}{1.3}
\begin{tabularx}{\linewidth}{@{}>{\bfseries}p{4.5cm} X@{}}
\toprule
\rowcolor{light} Métrique & Valeur observée \\ \midrule
Période d'échantillonnage  & 2 secondes \\
Débit messages             & $\sim$30 / minute \\
Taille payload JSON        & $\sim$80 octets \\
Latence FiPy $\rightarrow$ Grafana & $\sim$400 ms \\
Consommation FiPy seul     & 120 mA @ 3,3 V ($\sim$0,4 W) \\
PySense + capteurs         & +30 mA \\
LED RGB                    & +5 à 10 mA \\
\textbf{Total système}     & \textbf{$\sim$0,5 W} en régime normal \\
\bottomrule
\end{tabularx}

\medskip
La faible empreinte mémoire de MicroPython et l'efficacité du header MQTT (contre HTTP) permettent une exécution stable sur plusieurs heures sans fuite ni blocage.

\section{Difficultés rencontrées}

\renewcommand{\arraystretch}{1.4}
\begin{tabularx}{\linewidth}{@{}>{\bfseries}p{4cm} X@{}}
\toprule
\rowcolor{light} Problème & Solution \\ \midrule
Bornes Eduroam/RT-WiFi bloquent les ports MQTT & Passerelle smartphone en mode point d'accès \\
Incompatibilité de versions des libs Pycom & Mise à jour firmware FiPy et PySense avant déploiement \\
Mapping JSON $\leftrightarrow$ topics Grafana & Format de message uniformisé pour tous les capteurs \\
\bottomrule
\end{tabularx}

\section{Conclusion}

Ce TP matérialise la boucle IoT canonique : capter, transmettre, traiter, visualiser, agir. La maîtrise simultanée du firmware embarqué (MicroPython), du transport (MQTT/Wi-Fi), du stockage (MySQL) et de la visualisation (Grafana) constitue un socle solide pour des projets plus ambitieux \textemdash{} LoRa, LTE-M, architectures cloud distribuées.

L'intégration d'une commande MQTT inverse (bouton Grafana $\rightarrow$ LED FiPy) symbolise la boucle fermée : la donnée ne se contente plus d'être observée, elle devient un levier d'action.

\vspace{0.5cm}
\begin{center}
{\sffamily\itshape\color{midgray}\large
\og Internet of Things : connecter le monde physique au monde numérique. \fg{}}
\end{center}

\end{document}
